\documentclass{beamer}
\usetheme{Madrid}
\usecolortheme{default}

\usepackage{amsmath}
\usepackage{amssymb}
\usepackage{bm}
\usepackage{graphicx}
\usepackage{listings}
\usepackage{xcolor}

% Configure Python code formatting
\lstset{
    language=Python,
    basicstyle=\ttfamily\tiny,
    keywordstyle=\color{blue},
    stringstyle=\color{green!50!black},
    commentstyle=\color{red!80!black},
    showstringspaces=false,
    frame=single,
    breaklines=true
}

\title{Probability via Fundamental Axioms}
\subtitle{Problem 3.0.19 (MA 2007)}
\author{ee25btech11019}
\date{}

\begin{document}

\begin{frame}
    \titlepage
\end{frame}

\begin{frame}{Problem Statement}
    Let $A$, $B$, and $C$ be three events such that:
    \begin{equation*}
        P(A) = 0.4, \quad P(B) = 0.5, \quad P(A \cup B) = 0.6,
    \end{equation*}
    \begin{equation*}
        P(C) = 0.6, \quad \text{and} \quad P(A \cap B \cap C^c) = 0.1
    \end{equation*}
    
    Then $P(A \cap B \cap C) =$
    
    \vspace{0.5cm}
    a) $\frac{1}{2}$ \hspace{1cm} b) $\frac{1}{3}$ \hspace{1cm} c) $\frac{1}{4}$ \hspace{1cm} d) $\frac{1}{5}$
\end{frame}

\begin{frame}{Partitioning the Sample Space}
    The fundamental axiom that the sum of all mutually exclusive events in a sample space $\Omega$ must equal 1:
    \begin{equation}
        1 = P(A \cap B) + P(A \cap B^c) + P(A^c \cap B) + P(A^c \cap B^c)
    \end{equation}
    
    By the additivity axiom for disjoint sets, a marginal probability is the sum of its parts:
    \begin{align}
        P(A) &= P(A \cap B) + P(A \cap B^c) \\
        P(B) &= P(A \cap B) + P(A^c \cap B)
    \end{align}
\end{frame}

\begin{frame}{Substituting Marginals}
    We rearrange the marginal equations to isolate the exclusive regions:
    \begin{align}
        P(A \cap B^c) &= P(A) - P(A \cap B) \\
        P(A^c \cap B) &= P(B) - P(A \cap B)
    \end{align}
    
    Now, substitute these strictly into our sample space equation:
    \begin{equation*}
        1 = P(A \cap B) + [P(A) - P(A \cap B)] + [P(B) - P(A \cap B)] + P(A^c \cap B^c)
    \end{equation*}
    
    Simplifying the terms:
    \begin{equation}
        1 = P(A) + P(B) - P(A \cap B) + P(A^c \cap B^c)
    \end{equation}
\end{frame}

\begin{frame}{Applying De Morgan's Law}
    We now have the term $P(A^c \cap B^c)$. By De Morgan's Law, the probability of "neither A nor B" is the complement of their union:
    \begin{equation}
        P(A^c \cap B^c) = 1 - P(A \cup B)
    \end{equation}
    
    Substituting this back into our simplified equation:
    \begin{equation}
        1 = P(A) + P(B) - P(A \cap B) + 1 - P(A \cup B)
    \end{equation}
    
    Subtract 1 from both sides and isolate $P(A \cap B)$:
    \begin{equation}
        P(A \cap B) = P(A) + P(B) - P(A \cup B)
    \end{equation}
\end{frame}

\begin{frame}{Evaluating the Event Intersection}
    We have now derived the standard formula strictly from fundamental axioms without circular logic. We can substitute the given values:
    \begin{equation}
        P(A \cap B) = 0.4 + 0.5 - 0.6
    \end{equation}
    \begin{equation}
        P(A \cap B) = 0.3
    \end{equation}
    
    Finally, we partition this intersection using event $C$ (Law of Total Probability):
    \begin{align}
        P(A \cap B) &= P(A \cap B \cap C) + P(A \cap B \cap C^c) \\
        0.3 &= P(A \cap B \cap C) + 0.1 \\
        P(A \cap B \cap C) &= 0.2
    \end{align}
    
    \vspace{0.2cm}
    \textbf{Correct Option: (d) $\frac{1}{5}$}
\end{frame}

\begin{frame}[fragile]{Python Verification: Monte Carlo Simulation}
    To computationally verify our theoretical derivation, we simulate $10^6$ trials using the exact partitioned probabilities.
\begin{lstlisting}
import numpy as np

# Total number of simulations for the Law of Large Numbers
N = 1_000_000

# We partition the sample space into 5 strictly disjoint regions
# Region 1: A AND B AND C (Expected target = 0.2)
# Region 2: A AND B AND C_complement (Given = 0.1)
# Region 3: A AND not B (0.4 - 0.3 = 0.1)
# Region 4: B AND not A (0.5 - 0.3 = 0.2)
# Region 5: A^c \cap B^c (1 - P(A U B) = 0.4)

# Define the exact probabilities for these disjoint regions
probs = [0.2, 0.1, 0.1, 0.2, 0.4]

# Simulate drawing N samples from this probability distribution
# '0' represents our target intersection: A n B n C
events = np.random.choice([0, 1, 2, 3, 4], size=N, p=probs)

# Calculate the empirical probability
simulated_prob = np.sum(events == 0) / N

print(f"Theoretical P(A \u2229 B \u2229 C) = 0.2000")
print(f"Monte Carlo Simulated Probability (N={N}) = {simulated_prob:.4f}")
\end{lstlisting}
\end{frame}

\end{document}
